\section{Ergodic Theorem in Finite State Space}
\label{sec:A.4}

In this subsection, we study ergodicity of Markov chains under the assumption that the state space
$E$ is finite.  In this setting, irreducibility and aperiodicity are enough to guarantee ergodicity.
We will use the following lemma.

\begin{lemma}
\label{lem:A.22}
Let $P$ be the transition probability matrix of an irreducible, aperiodic, and finite state Markov
chain.  Then, there is an integer $m$ such that, for all $n \geq m$, the matrix $P^n$ has strictly
positive entries.
\end{lemma}

In order to establish ergodicity we will use a coupling argument which generalizes beyond the
finite state space setting considered here.  The coupling argument will be used repeatedly in these
notes, and we will see that central to the argument is bounding the Markov kernel from below,
which in the case of finite state space is achieved using the previous lemma.  The coupling
argument uses the total variation distance between probability measures.

\begin{definition}[Total Variation Distance]
\label{def:A.23}
The total variation distance between two probability measures $\nu_1, \nu_2$ on $E$ is defined by
\begin{equation}
  d_{\mathrm{TV}}(\nu_1, \nu_2) = \sup_{A \subseteq E} \abs{\nu_1(A) - \nu_2(A)}.
  \label{eq:A.1}
\end{equation}
\end{definition}

While \eqref{eq:A.1} is the standard definition of the total variation distance, we will use the
following characterization.

\begin{lemma}
\label{lem:A.24}
It holds that
\[
  d_{\mathrm{TV}}(\nu_1, \nu_2)
  = \frac{1}{2} \sup_{\abs{h}_\infty \leq 1}
    \bigl|\Expect_{X \sim \nu_1}[h(X)] - \Expect_{X \sim \nu_2}[h(X)]\bigr|.
\]
\end{lemma}

\begin{theorem}[Ergodicity of Markov Chains in Finite State Space]
\label{thm:A.25}
Let $\mathbb{X} = \{X_n\}_{n=0}^\infty$ be a Markov chain that is irreducible and aperiodic with
finite state space $E$.  Then, there are $\varepsilon > 0$ and a unique invariant distribution $\pi$
such that
\begin{equation}
  d_{\mathrm{TV}}(\pi_n, \pi) \leq (1 - \varepsilon)^n,
  \label{eq:A.2}
\end{equation}
where $\pi_n$ denotes the law of $X_n$.
\end{theorem}

\begin{proof}
We will assume (without loss of generality by Lemma~\ref{lem:A.22}) that $p_{ij} > \varepsilon$
for all $i, j \in E$.  Consider the following map from $\mathcal{P}(E)$ (the set of probability row
vectors on $E$) into itself:
\[
  \mathcal{P}(E) \to \mathcal{P}(E), \qquad \nu \mapsto \nu P.
\]
Since this map is continuous and $\mathcal{P}(E)$ is compact and convex, Brouwer's fixed point
theorem guarantees the existence of a fixed point.  That is, there is $\pi \in \mathcal{P}(E)$
such that $\pi P = \pi$, showing the existence of an invariant distribution.  In the rest of the
proof we will prove that if $\pi$ is an invariant distribution, then inequality~\eqref{eq:A.2}
holds.  This in particular will automatically imply the uniqueness of the invariant distribution.

We will use a \emph{coupling} argument.  Let $\{B_n\}_{n=0}^\infty$ be a sequence of independent
$\mathrm{Bernoulli}(\varepsilon)$ random variables, independent of all other randomness, and define
for given random variable $W_0$ and $n \geq 0$
\begin{equation}
  W_{n+1} \sim
  \begin{cases}
    s(W_n, \cdot) & \text{if } B_n = 0, \\
    r(W_n, \cdot) & \text{if } B_n = 1,
  \end{cases}
  \label{eq:A.3}
\end{equation}
where
\[
  s(i, j) = \frac{p_{ij} - \varepsilon r(i, j)}{1 - \varepsilon}
\]
and $r$ is the uniform transition kernel, so that $r(i, \cdot)$ is the discrete uniform
distribution in $E$.  Note that
\[
  s(i, E) = \frac{p(i, E) - \varepsilon r(i, E)}{1 - \varepsilon}
           = \frac{1 - \varepsilon}{1 - \varepsilon} = 1
\]
and
\[
  s(i, j) \geq \frac{\varepsilon - \varepsilon r(i, j)}{1 - \varepsilon} \geq 0,
\]
so $s$ is a Markov kernel.  Moreover,
\begin{align*}
  \Prob(W_{n+1} = j \mid W_n = i)
    &= \varepsilon \Prob(W_{n+1} = j \mid B_n = 1, W_n = i)
       + (1 - \varepsilon)\Prob(W_{n+1} = j \mid B_n = 0, W_n = i) \\
    &= \varepsilon r(i, j) + p(i, j) - \varepsilon r(i, j) \\
    &= p(i, j).
\end{align*}
Thus, $\{W_n\}$ has transition kernel $P$.

Let $h : E \to \R$ with $\abs{h}_\infty \leq 1$, and let $\tau = \min(n \in \N : B_n = 1)$.
Regardless of the distribution of $W_0$,
\begin{align*}
  \Expect[h(W_n)]
    &= \Expect[h(W_n) \mid \tau \geq n]\,\Prob(\tau \geq n)
       + \sum_{l=0}^{n-1} \Expect[h(W_n) \mid \tau = l]\,\Prob(\tau = l) \\
    &= \underbrace{\Expect[h(W_n) \mid \tau \geq n]}_{\leq 1}
       \underbrace{\Prob(\tau \geq n)}_{\leq (1-\varepsilon)^n}
       + \sum_{l=0}^{n-1}
         \underbrace{\Expect_{W_0 \sim \mathrm{Unif}(E)}[h(W_{n-l})]}_{\text{independent of the initial distribution}}
         \Prob(\tau = l),
\end{align*}
where $\mathrm{Unif}(E)$ denotes the uniform distribution on $E$.

Now define two sequences of random variables as in equation~\eqref{eq:A.3}: let $\{Y_n\}_{n=0}^\infty$
be initialized with $\pi_0$ (the initial distribution of $\mathbb{X}$), and let $\{Z_n\}_{n=0}^\infty$
be initialized with an invariant distribution $\pi$ (which we have shown that exists) so that
$Z_n \sim \pi$ for all $n$.  Then,
\[
  d_{\mathrm{TV}}(\pi_n, \pi)
  = \frac{1}{2} \sup_{\abs{h}_\infty \leq 1}
    \Bigl|\Expect_{Y_n \sim \pi_n}[h(Y_n)] - \Expect_{Z_n \sim \pi}[h(Z_n)]\Bigr|
  \leq (1 - \varepsilon)^n. \qedhere
\]
\end{proof}
